

Deploying synthetic clinical data hinges on a trade-off between preserving the non-linear dependencies of clinical variables and limiting re-identification risk. Our 23-benchmark evaluation suggests hierarchical reinforcement learning navigates this trade-off while revealing design boundaries that motivate future work.

\textbf{Structural and Semantic Consistency through Hierarchical Feedback.}
A primary failure mode of generative models in healthcare is producing implausible records with contradictory attributes. Adversarial models such as CTGAN learn these constraints only implicitly and often fail in high-dimensional discrete spaces (vanishing gradients, mode collapse). DISCO-TAB instead decomposes record consistency into explicit levels of verification, separating token-level syntactic consistency from row-level semantic and logical coherence. Automated Constraint Discovery reinforces this by identifying statistically significant feature dependencies (e.g., between body-mass index and diabetes progression), which are well preserved in the output (mean Pearson correlation above $0.88$ on every dataset), biasing the generator toward medically plausible regions that unguided LLM approaches lack. This discovery is currently pairwise, so DISCO-TAB enforces first-order and pairwise clinical logic rather than full causal reasoning; causal discovery or clinician-in-the-loop validation could address higher-order dependencies.

\textbf{Robustness in Data-Scarce Regimes.}
Precision medicine targets rare phenotypes with inherently limited data ($N<500$), where likelihood-based Transformers and diffusion models often overfit or inject excessive noise, driving utility toward chance or collapse. DISCO-TAB is among the most reliable generators we evaluate: top-tier utility among full-coverage methods, leading on correlation preservation, and no single-class collapse (unlike TabSyn, GReaT, and TVAE on the smallest, most imbalanced sets). Performance on the tiniest sets does vary across seeds, so we report multi-seed means throughout. This robustness arises from reframing generation as reward maximization rather than likelihood estimation: multi-granular reinforcement signals steer the generator toward valid regions of sparse clinical manifolds. The cost is training overhead, several times a one-shot GAN's, from the discriminators plus PPO; inference stays efficient. This is a deliberate trade of training efficiency for stability and fidelity in scarce-data regimes, and parameter-efficient fine-tuning, adaptive discriminator activation, or lighter RL objectives could reduce it.

\textbf{Navigating the Privacy--Utility Frontier.}
Diffusion models such as TabDDPM reach large distances to real records by injecting noise that displaces samples off the manifold, raising apparent privacy while degrading utility, a privacy--utility collapse rather than a desirable regime. DISCO-TAB instead holds a pragmatic middle ground, preserving the structure needed for high-utility learning, which suits collaborative and federated settings where trust depends on both. It also inherits its pre-trained backbone's biases: the row-level discriminator mitigates many hallucinations, but domain-adaptive or medically grounded pretraining would further curb incorrect priors.

\rev{\textbf{Practical Implications.}
DISCO-TAB targets workflows where access to clinical tabular data is constrained: prototyping and pre-screening machine-learning pipelines before touching protected data, augmenting under-represented classes, and benchmarking. Because it offers no formal privacy guarantee, its records should \emph{not} be assumed safe for unrestricted release; deployment should follow a dataset-specific disclosure-risk assessment with institutional governance, and a differentially private mechanism where formal guarantees are required.}

\rev{\textbf{Advantages and Limitations.}
The principal advantages of DISCO-TAB are its ability to combine heterogeneous validity signals within a single RL objective, its explicit attention to minority-class preservation, its automated discovery of statistically salient dependencies, and its applicability across heterogeneous datasets. These advantages are counterbalanced by several concrete limitations, which we state explicitly. (i) Disclosure risk is characterized empirically, through DCR and related proximity measures, rather than by a formal differential-privacy guarantee. (ii) Automated constraint discovery is currently pairwise and cannot capture higher-order or causal interactions. (iii) The reward is a terminal, sequence-level signal aggregated from multiple granularities rather than a temporally dense, per-step signal, so credit assignment across generation steps is coarse. (iv) Reward quality depends on the calibration of the learned discriminators, which may themselves be biased or exploitable (reward hacking), and performance is sensitive to the reward weights $\lambda_k$, the IFRS factor $\alpha$, and the serialization order. (v) The inverse-frequency reward shaping depends on the choice of target label. (vi) Training incurs the computational overhead of PPO plus several discriminators. (vii) Evaluation is retrospective and public-benchmark-based; we did not perform clinician review of clinical plausibility or validate on large multi-institutional EHR cohorts.}

\textbf{Outlook.}
These findings position DISCO-TAB as a principled framework for high-fidelity biomedical tabular synthesis, with strong structural and semantic consistency, robustness in data-scarce regimes, and a balanced privacy--utility trade-off. Future work will reduce computational overhead, extend constraint discovery beyond pairwise statistics, compare alternative policy-optimization objectives such as Direct Preference Optimization (DPO) against the PPO loop used here, and incorporate formal privacy guarantees.





